% The template needs to be compiled using LuaLaTeX or XeLaTeX

\documentclass{article}
\usepackage{jtdh}
% \definecolor{Mycolor2}{HTML}{00F9DE}
\usepackage[utf8]{inputenc} % allow utf-8 input
\usepackage[slovene]{babel} % Change to english if english paper
\usepackage{url}   % simple URL typesetting
\AtBeginDocument{\urlstyle{APACsame}}
\usepackage[hang,flushmargin]{footmisc} % no indentation in footnotes
%\usepackage[colorlinks = true, linkcolor = black, urlcolor  = gray, citecolor = black, anchorcolor = black]{hyperref}
\usepackage{booktabs}       % professional-quality tables
\usepackage{amsfonts}       % blackboard math symbols (already provided via unicode-math in jtdh.sty)
%\usepackage{amssymb}        % dodatni simboli (conflicts with unicode-math under LuaLaTeX)
\usepackage{amsmath}        % eqref, etc. (loaded transitively via unicode-math)
\usepackage{nicefrac}       % compact symbols for 1/2, etc.
\usepackage{xcolor}
\usepackage{graphicx}
\usepackage{tabularx}
\usepackage{longtable}      % Za dolge tabele
\usepackage[]{apacite}
\usepackage{enumitem}
\usepackage{bm}
% Improve float and pagination control
\usepackage{placeins} % provides \FloatBarrier to keep floats from crossing logical boundaries
% Reduce risk of overfull lines without changing content
\emergencystretch=3em
% Geometry warning fix: ensure margin notes don't overrun the paper (even if we don't use them)
\setlength{\marginparwidth}{40pt}
\setlength{\marginparsep}{6pt}
\graphicspath{{Images/}{img/}} % Prilagojeno iz diplomske, original je bil 'img/'
% Improve table spacing: avoid vertically compressed tables and legend overlap
% Increase caption separation and float spacings conservatively
\setlength{\abovecaptionskip}{6pt}
\setlength{\belowcaptionskip}{8pt}
\setlength{\textfloatsep}{12pt plus 2pt minus 2pt}
\setlength{\floatsep}{10pt plus 2pt minus 2pt}
\setlength{\intextsep}{10pt plus 2pt minus 2pt}
% Default table row spacing (increase slightly). Local tables may override this.
%\renewcommand{\arraystretch}{1.05}
% table styles
\aboverulesep=0ex
\belowrulesep=0ex
\raggedbottom


%\addbibresource{biblio.bib} %Imports bibliography file 

\title{Odstranjevanje mašil v govoru}

\author{
  Urban Vesel\textsuperscript{1}, 
}
\institutions{
  \textsuperscript{1}Univerza v Ljubljani, Fakulteta za računalništvo in informatiko
}

\begin{document}
\maketitle

\begin{abstract}

\end{abstract}
\keywords{}


\section{Uvod}



\bibliographystyle{apacite}
\urlstyle{same}
\bibliography{biblio}


\newpage
\titleeng{Removal of filler words from speech}
\begin{abstracteng}
Fillers are words, phrases, or sounds that do not contribute to the meaning of speech and are often distracting. 
This paper addresses the problem of automatic filler removal in Slovenian speech. 
The main basis is provided by modern automatic speech recognition (ASR) models based on neural network architectures. 
These are end-to-end models that, using deep learning on large amounts of audio recordings and corresponding transcriptions, can recognize speech. 
Depending on the training data and model architecture, certain models support word-level time stamps. 
This capability is already used in some foreign languages for detecting and subsequently removing fillers in speech, while no such tool yet exists for Slovenian.
Within the scope of this thesis, 
we developed software tools that enable easy use and comparison of different ASR models for detecting and removing fillers in Slovenian audio recordings. 
The results show that ASR models, with appropriate adaptation, offer a promising solution for automatic filler removal in Slovenian speech, 
while the developed tools open possibilities for further research and improvements in the field of speech technology.
\end{abstracteng}

\keywordseng{automatic speech recognition (ASR), 
automatic filler word removal, filler word detection, Slovenian language}

\end{document}